% Abstract (TCC2, results-aware; leads with both aims + the Phi/P_LLM dissociation).
% Single paragraph; keywords set in main.tex via \keywords{}.

This work pursues two pre-registered aims: (i) to formalize quantitative metrics for detecting RLHF-induced behavioral distortion and for measuring behavioral fidelity in LLM-based social simulations, and (ii) to use those metrics to test whether empirical grounding in real socioeconomic microdata propagates through instruction-tuned LLM agents into realistic collective behavior. Instruction-tuned LLMs deployed as agents in multi-agent economic simulations exhibit a systematic anomaly --- cooperation rates exceeding both the Nash equilibrium and laboratory human behavior --- which we formalize as the \textbf{RLHF Cooperative Bias} and operationalize as the total-variation distance of the observed action distribution from the uniform action prior, B\_RLHF = TV(π, π\_uniform); behavioral fidelity is measured by a composite \textbf{Behavioral Realism Metric (BRM ∈ [0,1])}. Both metrics are embedded in the \textbf{Behavioral Grounding Framework}, a formally specified platform BGF = (A, E, G, P, Φ, T) grounded in European Social Survey Round 11 microdata via a population-synthesis map Φ: D\_ESS → Profile. The central finding is a \textbf{dissociation (Φ/P\_LLM)}: empirical grounding that is effective at the rule-based policy layer does \textbf{not} propagate through the LLM decision channel. The pre-registered 10-seed N=100 test finds the grounded and ungrounded arms statistically indistinguishable on cooperation (0.461 vs 0.455, MWU p = 0.91), Gini, and wealth --- the pre-registered grounding hypothesis is \textbf{not supported at primary scale}, and we report this negative result openly. The metrics are precisely what surface the dissociation: at the layer that reads grounding directly, a deterministic policy reproduces a European-range Gini coefficient (0.325, BCa 95\% CI [0.324, 0.325]; N=500, T=30, 10 seeds) and the cross-cultural cooperation gradient (Spearman ρ = +1.000; World Values Survey r = +0.977; Herrmann--Thöni--Gächter public-goods benchmark ρ = +0.886, p = 0.033), establishing that Φ is empirically valid while the LLM decision layer remains RLHF-anchored. A scale-dependent RLHF cooperation cascade (B\_RLHF amplified 2.6--3.2× at N=500) and a falsified memory-depth hypothesis corroborate this anchoring reading. The framework, the two metrics, and the dissociation are released as an open-source artifact (1{,}578 automated tests, one-command reproduction, cryptographic reproducibility witness) so the result can be re-tested in any LLM-as-agent setting on any empirical population.
